% Auto-generated by generate_tex_fragments.py -- do not hand-edit.
\begin{longtable}{@{}L{2.0cm}L{2.0cm}L{2.1cm}L{6.2cm}L{1.9cm}@{}}
\caption{P5: Referral quality depends on timeliness, accessibility and follow-through, not offer alone. Record, design, population/setting, finding as charted (verbatim from the library record), and direction relative to the proposition. n=37 linked records (kg/graph.json edges).}\label{tab:p5}\\
\toprule
\textbf{Record} & \textbf{Design} & \textbf{Population\slash{}setting} & \textbf{Finding (as charted)} & \textbf{Direction}\\
\midrule
\endfirsthead
\multicolumn{5}{l}{\textit{P5 continued}}\\
\toprule
\textbf{Record} & \textbf{Design} & \textbf{Population\slash{}setting} & \textbf{Finding (as charted)} & \textbf{Direction}\\
\midrule
\endhead
\bottomrule
\endfoot
Alghamdi M 2022 \cite{srp34253076} & \seqsplit{Qualitative} dominant design with \seqsplit{supplemental} \seqsplit{quantitative} data (n=8) & Canadian Muslim immigrant women, 6 countries of origin & Pre-migration adverse childhood \seqsplit{experiences} and family history of IPV, combined with post-migration language\slash{}social-isolation barriers and \seqsplit{internalized} \seqsplit{patriarchal} values, shaped help-seeking and the eventual decision to leave abusive \seqsplit{relationships} among 8 Muslim immigrant women. & Supports\\
\\[3pt]
Azalee M 2025 \cite{srp40116416} & \seqsplit{Qualitative} in-depth interviews (n=6) & Health care workers who were IPV survivors, tertiary hospital, Malaysia & Six HCW-survivors used varied coping strategies (seeking help, \seqsplit{spirituality}, avoidance, self-harm, staying\slash{}leaving); survivors relied mainly on coworkers and sought formal help only when situations became critical, hindered by \seqsplit{misconceptions} about IPV, privacy concerns and fear of workplace gossip. & Supports\\
\\[3pt]
Baker CK 2022 \cite{srp34838218} & Systematic review of 34 controlled \seqsplit{quantitative} studies (from 23,902 screened records) & Women \seqsplit{experiencing} IPV, \seqsplit{predominantly} US-based housing\slash{}shelter programmes & [CORRECTION: first\_author is wrong: search lists 'Baker CK \slash{} Sprague S (co-first authorship per abstract group)' but PubMed indexes the first author as Yakubovich (no Baker or Sprague in the author-derived first-author field). Content confirmed: 'screened 23 902 unique records and identified 34 eligible studies'; 'no cumulative evidence of \seqsplit{disadvantages} following any IPV-housing \seqsplit{intervention}. Evidence of benefits was strongest for mental health outcomes, intent to leave partner, perceived safety, and housing and partner-related stress'; 'more research...needed...from outside of the USA.' FIX: change first\_author to Yakubovich.] No cumulative evidence of \seqsplit{disadvantage} from any IPV-housing \seqsplit{intervention}; evidence of benefit was strongest for mental health outcomes, intent to leave partner, perceived safety, and reduced housing\slash{}partner-related stress; most included studies were at high risk of bias, and evidence for long-term\slash{}non-US housing solutions was sparse. & Supports\\
\\[3pt]
Chen 2013 \cite{ref94} & clinical practice review & primary\slash{}family care settings & Describes RADAR (Remember to ask routinely, Ask directly in private, Document, Assess safety, Review options) as a structured clinical tool for \seqsplit{identifying} and responding to IPV, \seqsplit{emphasizing} private \seqsplit{questioning}, safety assessment (weapons, children, \seqsplit{perpetrator} risk), and advocacy-based referral requiring trained links to community resources. & Supports\\
\\[3pt]
Feltner 2018 \cite{srp30357304} & systematic review \seqsplit{commissioned} for USPSTF (global-level guidance body), 30 included studies & adults in healthcare settings, US and comparable settings & [CORRECTION: PMID\slash{}title\slash{}year correct, but the key\_finding materially overstates the abstract's \seqsplit{conclusions} in two ways. (1) The abstract's actual conclusion is: 'trials of IPV screening in adult women did NOT show a reduction in IPV or \seqsplit{improvement} in quality of life over 3 to 18 months' --- screening tools were found to 'reasonably identify' IPV (accuracy claim only, \seqsplit{sensitivity}\slash{}\seqsplit{specificity} varied widely by tool), but the record's claim that the review found 'some \seqsplit{interventions} (e.g., counseling) following screening among women of \seqsplit{reproductive} age reduce IPV' is too broad: only 2 of 11 \seqsplit{intervention} RCTs (1 home-visiting, 1 multi-risk behavioral counseling) showed \seqsplit{significant} reduction, and only in PREGNANT women \seqsplit{specifically}, not the general \seqsplit{reproductive}-age population. (2) 'informing a USPSTF grade-B \seqsplit{recommendation}' conflates this article (the evidence report\slash{}systematic review, PMID 30357304) with the companion USPSTF \seqsplit{Recommendation} Statement, which is a SEPARATE PMID (30357305, also JAMA 2018;320(16):1678-1687) not cited in the record. Fix: rewrite key\_finding to state the evidence report found screening tools moderately accurate but screening ITSELF did not reduce IPV\slash{}improve QoL in trials; only 2 specific \seqsplit{interventions} in pregnant-women subgroups showed benefit; and either cite PMID 30357305 separately for the actual grade-B \seqsplit{recommendation} or explicitly note this record is the evidence report underlying that \seqsplit{recommendation}, not the \seqsplit{recommendation} itself.] Concludes screening \seqsplit{instruments} can accurately identify IPV in health care settings and that some \seqsplit{interventions} (e.g., counseling) following screening among women of \seqsplit{reproductive} age reduce IPV; harms of screening were assessed and found to be low, informing a USPSTF grade-B \seqsplit{recommendation} for \seqsplit{reproductive}-age women. & Supports\\
\\[3pt]
Garcia-Moreno C 2015 \cite{srp25467583} & Series paper \slash{} narrative review with five country case studies & Health systems in low-, middle- and high-income countries & Reviews evidence for clinical \seqsplit{interventions} and health-system components (protocols, training, capacity-building) needed to identify and support women subjected to IPV\slash{}sexual violence following the 2013 WHO guidelines and 67th World Health Assembly resolution; finds \seqsplit{substantial} system and \seqsplit{behavioural} barriers to \seqsplit{implementation}, especially in low- and middle-income countries, and that system \seqsplit{development} has generally been slow. & Supports\\
\\[3pt]
Gerth J 2014 \cite{ref168} & Narrative\slash{}validation review & IPV cases assessed by police, hospital and victim-services staff (multi-\seqsplit{professional}) & [CORRECTION: Actual PubMed title: '[The Ontario Domestic Assault Risk Assessment (ODARA) - Validity und authorised German \seqsplit{translation} of an intimate partner violence screening tool].' Record has 'Validity and authorised German \seqsplit{translation}' --- wrong \seqsplit{conjunction} ('and' vs the German 'und', which PubMed retains in the bilingual title) and drops the subtitle 'of an intimate partner violence screening tool'. FIX: correct the title string. Key finding is accurate: abstract quote 'provide a brief risk assessment tool which can be scored on the basis of only few and easily \seqsplit{collectable} \seqsplit{information}... wider range of \seqsplit{practitioners} (e.g., police officers, hospital and victim services' staff).'] ODARA was designed to be brief and scorable by a wide range of \seqsplit{practitioners} (police, hospital, victim services) using easily \seqsplit{collectable} \seqsplit{information}, to reliably \seqsplit{discriminate} low- from high-risk cases and guide \seqsplit{intervention}. & Supports\\
\\[3pt]
Hill AL 2019 \cite{srp31306331} & Secondary cross-sectional analysis of baseline survey data & 550 sexually active female high school students, USA (school health center attendees, ages 14-19) & [CORRECTION: Prevalence numbers are correct (12\% recent RC, 17\% ARA) among 550 sexually active female students. PROBLEM: key\_finding claims 'both were associated with \seqsplit{differences} in care-seeking and sexual health behaviors' -- but the abstract EXPLICITLY states the opposite for \seqsplit{reproductive} coercion alone: 'There were no observed \seqsplit{significant} \seqsplit{differences} in care-seeking behaviors among those with recent \seqsplit{reproductive} coercion compared with those without.' It is adolescent \seqsplit{relationship} abuse (ARA), not \seqsplit{reproductive} coercion, that was associated with a care-seeking difference (increased odds of STI testing\slash{}treatment, aOR 2.08); the combined ARA+RC-exposed group showed \seqsplit{differences} in sexual-partner\slash{}\seqsplit{contraception} behaviors, not general 'care-seeking'. FIX: rewrite key\_finding to state that RC alone showed NO \seqsplit{significant} care-seeking difference, ARA alone was associated with increased STI-testing care-seeking (aOR 2.08), and the ARA+RC combined group showed elevated odds of specific risk behaviors (older partner, more partners, hormonal-only \seqsplit{contraception}) -- do not state both constructs were 'associated with \seqsplit{differences} in care-seeking'.] 12\% reported recent \seqsplit{reproductive} coercion and 17\% reported adolescent \seqsplit{relationship} abuse; both were associated with \seqsplit{differences} in care-seeking and sexual health behaviors, measured via validated RC and \seqsplit{relationship}-abuse \seqsplit{instruments}. & Supports\\
\\[3pt]
Le V 2025 \cite{srp37649394} & \seqsplit{Qualitative} \seqsplit{implementation} study (16 semi-structured interviews, CFIR-coded) & Catholic diocesan and parish leadership, \seqsplit{Archdiocese} of Chicago Domestic Violence Outreach, USA & [CORRECTION: Key\_finding content (seven CFIR inner-setting constructs and their content) is supported by the abstract. However, first\_author is wrong: record lists 'Le V (\seqsplit{Archdiocese} of Chicago Domestic Violence Outreach study)' but PubMed's actual first author is 'Debinski B'. FIX the first\_author field.] Using the \seqsplit{Consolidated} Framework for \seqsplit{Implementation} Research, seven inner-setting constructs (leadership structure, culture, resources, networking\slash{}\seqsplit{communication}, access to knowledge, \seqsplit{compatibility} with parish values, leadership engagement) were identified as \seqsplit{facilitators} of successful diocesan-level IPV response \seqsplit{implementation}. & Supports\\
\\[3pt]
Lee 2019 \cite{srp31012540} & pre-post quality \seqsplit{improvement} \seqsplit{intervention} study & healthcare providers (nurses, physicians, social workers) at a US clinical site & An EMR-integrated brief IPV screening tool plus provider education plus an automated resource\slash{}referral telephone system was associated with measured \seqsplit{improvement} in provider readiness to screen (assessed via the Domestic Violence Health Care Provider Survey Scale, pre\slash{}post t-test). & Supports\\
\\[3pt]
McCauley HL 2017 \cite{ref17} & Secondary \seqsplit{psychometric} analysis of pooled \seqsplit{multicenter}\slash{}RCT survey data & 4,674 women aged 16-29 in 24 \seqsplit{Pennsylvania} and 5 California family planning clinics, USA & Confirmed a 2-domain structure (pregnancy coercion, condom \seqsplit{manipulation}) and produced a validated 5-item short-form \seqsplit{Reproductive} Coercion Scale (RCS); recent \seqsplit{reproductive} coercion prevalence was 6.3-6.7\% depending on the scale version used. & Supports\\
\\[3pt]
Meiksin R 2020 \cite{srp32460786} & Trial protocol & Women\slash{}girls seeking \seqsplit{contraceptive} services, community-based clinics, Nairobi, Kenya & [CORRECTION: first\_author is wrong: search lists 'Meiksin R' but PubMed indexes the first author as Uysal. Content confirmed: 'ARCHES is a single-session, clinic-based model...Core elements...\seqsplit{opportunity} for patient disclosure of RC and IPV (and subsequent warm referral to local services)' and 'To date, clinic-based GBV \seqsplit{interventions} have not been shown to reduce RC or unintended pregnancy in LMIC settings.' FIX: change first\_author to Uysal.] Describes ARCHES as a single-session, clinic-based enhanced counseling model that creates \seqsplit{opportunity} for patient disclosure of \seqsplit{reproductive} coercion\slash{}IPV followed by warm referral to local services; explicitly notes that to date, no clinic-based GBV \seqsplit{intervention} had been shown to reduce \seqsplit{reproductive} coercion or unintended pregnancy in LMIC settings, motivating this trial. & Supports\\
\\[3pt]
Messing JT 2013 \cite{ref164} & Meta-analytic review of predictive validity (AUC) & Five IPV risk-assessment \seqsplit{instruments} (ODARA, SARA, DA, DVSI, K-SID) across multiple validation studies & ODARA had highest weighted AUC=.666 (k=5), then SARA AUC=.628, DA AUC=.618, DVSI AUC=.582, K-SID AUC=.537; only 9\slash{}20 (45\%) of predictive-validity estimates involved the instrument \seqsplit{administered} correctly per its intended protocol. & Supports\\
\\[3pt]
Nagae M 2004 \cite{ref105} & \seqsplit{Qualitative} interview study & Abused women and health\slash{}legal\slash{}community actors, Khon Kaen province, \seqsplit{northeastern} Thailand & Interviews with 4 abused women and \seqsplit{surrounding} health \seqsplit{professionals}, government\slash{}legal staff and community members in Khon Kaen showed health \seqsplit{professionals} (nurses, community health workers) and, in one case, a community religious leader (a Buddhist monk) played \seqsplit{identifiable} roles in \seqsplit{identifying}, treating and supporting abused women within a multi-actor referral network. & Supports\\
\\[3pt]
O'Doherty L 2015 \cite{ref38} & Cochrane systematic review and meta-analysis of 13 RCTs\slash{}quasi-RCTs (n=14,959) & Women in antenatal, women's health, emergency department and primary care settings, mostly high-income urban settings & Screening increased clinical \seqsplit{identification} of IPV victims (OR 2.95, 95\% CI 1.79-4.87, n=10,074) but there was no evidence of effect on referrals to support services (OR 2.24, 95\% CI 0.64-7.86, only 2 studies, n=1298), no evidence screening reduced women's re-exposure to violence at 3-18 months, and \seqsplit{insufficient} evidence on health outcomes or harm beyond the immediate post-screening period; authors conclude screening alone is \seqsplit{insufficiently} supported to justify universal \seqsplit{implementation} absent stronger referral\slash{}outcome evidence. & Supports\\
\\[3pt]
Sharifnia A 2025 \cite{srp39425490} & Meta-\seqsplit{ethnography} (33 \seqsplit{qualitative} studies \seqsplit{synthesized}), 6 databases searched & Muslim women, global (multi-country \seqsplit{qualitative} evidence) & Synthesis of 33 \seqsplit{qualitative} studies identified four major themes - experience\slash{}impact of abuse, risk factors, help-seeking, and coping - and found religion functions both as a risk factor for abuse and a barrier to help-seeking, and \seqsplit{simultaneously} as a coping resource for some women. & Supports\\
\\[3pt]
Sprague 2017 \cite{srp28649297} & scoping review, 43 included studies & health care settings, multiple countries & Of 43 studies across 9 categories of IPV assistance programmes (\seqsplit{counselling}\slash{}advocacy, safety planning, referral, resource provision, home visitation, case management, videos, provider cueing, system changes), programmes using a single active-assistance type combined with a counsellor\slash{}advocate reported positive results in 77.8\% of studies, suggesting simpler, advocate-led \seqsplit{interventions} outperform complex bundled ones. & Supports\\
\\[3pt]
Trabold N 2014 \cite{ref109} & 3-year community-based \seqsplit{participatory} mixed-methods \seqsplit{demonstration} project across 4 community health centers & Urban, \seqsplit{predominantly} indigent safety-net primary care population, USA & [CORRECTION: first\_author is wrong: search lists 'Trabold N' but PubMed indexes the first author as Rhodes. Design \seqsplit{description} matches ('3-year community-based mixed-method \seqsplit{participatory} research project involving four community health centers...\seqsplit{predominately} indigent urban population'). key\_finding's general claim of '\seqsplit{significant} system-level challenges' is supported at a high level ('Results highlight the challenges, but also the \seqsplit{opportunities}, for \seqsplit{implementing} the new USPSTF guidelines...in resource-limited primary care settings') but the abstract itself does not itemize 'workflow, staff capacity, referral \seqsplit{infrastructure}' as specific named barriers -- that level of detail is not verifiable from the abstract alone (may be in full text). FIX: change first\_author to Rhodes; consider softening the \seqsplit{parenthetical} specifics or citing they come from the full text, not the abstract.] \seqsplit{Implementing} USPSTF-guideline-consistent universal IPV screening and referral in resource-limited urban primary care settings faced \seqsplit{significant} system-level challenges (workflow, staff capacity, referral \seqsplit{infrastructure}), even with a dedicated multi-year system-level \seqsplit{intervention} designed to raise detection and response. & Supports\\
\\[3pt]
US Preventive Services Task Force 2025 \cite{ref40} & guideline \slash{} \seqsplit{recommendation} statement (systematic-review-based) & US \seqsplit{adolescents} and adults, pregnant\slash{}postpartum and \seqsplit{reproductive}-age women; older\slash{}vulnerable adults without recognized signs of abuse & USPSTF concludes screening for IPV in women of \seqsplit{reproductive} age (including pregnant\slash{}postpartum), combined with provision or referral to \seqsplit{multicomponent} \seqsplit{interventions}, has moderate net benefit (B \seqsplit{recommendation}). For caregiver abuse\slash{}neglect of older or vulnerable adults, benefits and harms remain \seqsplit{undetermined} (I statement) even in this 2025 update. & Supports\\
\\[3pt]
Waalen J(CFPSA 2014 \cite{ref111} & Systematic review of 9 RCTs of physician-directed domestic violence education & \seqsplit{Postgraduate} trainee and practising physicians & [CORRECTION: first\_author is wrong: search lists 'Waalen J' but PubMed indexes the first author as Zaher. Content confirmed precisely: 'Brief \seqsplit{interventions} for \seqsplit{postgraduate} trainee physicians improved knowledge but did not seem to affect behaviour...Physician training combined with system support \seqsplit{interventions} seemed to benefit domestic violence victims and increase referrals to domestic violence support resources.' FIX: change first\_author to Zaher.] Brief \seqsplit{educational} \seqsplit{interventions} for \seqsplit{postgraduate} trainees improved knowledge but did NOT change screening\slash{}\seqsplit{identification} behaviour; only \seqsplit{multifaceted} training combined with system-support \seqsplit{interventions} (reminders, access to victim-support referral pathways, audit\slash{}feedback) was associated with increased referrals to domestic-violence support resources. & Supports\\
\\[3pt]
Wang PL 2015 \cite{ref165} & Validation cohort study & 543 female IPV victims in a Taiwanese victim-services program & TIPVDA showed strong \seqsplit{discriminant} power and ROC-AUC support for predicting lethal-danger risk, with stronger predictive power at the high-\seqsplit{dangerousness} end. & Supports\\
\\[3pt]
Williamson HC 2014 \cite{ref106} & Cross-sectional survey study (N=2,126 married \seqsplit{individuals}) & US married adults, general population sample & Contrary to the assumption that premarital education reduces later need for counseling, receiving premarital education covaried with an INCREASED likelihood of later couples counseling (a 'gateway' effect), with the \seqsplit{association} stronger among African Americans and those with lower income\slash{}education. & Supports\\
\\[3pt]
Williamson HC 2019 \cite{ref107} & Cross-sectional self-report study (N=231 ethnically diverse newlywed couples in low-income \seqsplit{neighborhoods}) & US low-income newlywed couples & Top perceived barriers to seeking couple therapy were cost and \seqsplit{uncertainty} about where to go for help; direct experience with premarital education (or indirect experience via a social-network member's couple therapy) was associated with reduced barriers and greater likelihood of \seqsplit{subsequently} seeking therapy. & Supports\\
\\[3pt]
Zust BL 2021 \cite{ref118} & Repeated cross-sectional survey (\seqsplit{convenience} sample of Upper-Midwest US \seqsplit{congregations}, 2005 and 2015) & Christian clergy and \seqsplit{congregation} members, US & [CORRECTION: Key\_finding is well supported: 'Victims find strength in their faith and would rather endure the violence at all costs to keep a family or a marriage together, than to compromise their faith by leaving' and 'Clergy felt \seqsplit{comfortable} in making referrals for \seqsplit{professional} counseling.' However relevance\_to\_tphe overstates the referral point: the abstract's very next clause is 'while the majority of members would prefer counseling with their pastor if they were in a violent \seqsplit{relationship}' --- i.e. \seqsplit{congregants} actually want direct pastoral counseling, not referral, and clergy said they want MORE training to counsel\slash{}speak about DV themselves, not that they prefer referral over direct counseling. FIX: change 'clergy themselves report needing referral pathways rather than attempting DV counseling directly' to something like 'clergy report comfort in making referrals, while members\slash{}clergy \seqsplit{simultaneously} want more direct pastoral counseling capacity --- a live tension rather than clean support for a refer-not-counsel policy that T-PHE's safeguard should note, not paper over.'] Faith \seqsplit{commitments} led some domestic-violence victims to endure abuse rather than compromise their faith by leaving; over 10 years, \seqsplit{congregational} support for victims changed slowly; clergy wanted more training in counseling and speaking about domestic violence, and felt \seqsplit{comfortable} making referrals but not \seqsplit{necessarily} providing direct DV counseling themselves. & Supports\\
\\[3pt]
\seqsplit{Grisurapong} S 2002 \cite{ref104} & \seqsplit{Intervention}\slash{}comparison case-site \seqsplit{description} & Women \seqsplit{experiencing} physical or sexual assault, Khon Kaen provincial hospital, Thailand & Describes the \seqsplit{establishment} and early operation of a hospital-based One Stop Crisis Center (OSCC) for women survivors of violence in Khon Kaen, Thailand, compared against a non-\seqsplit{intervention} provincial hospital site. & Supports\\
\\[3pt]
Hogan NR 2024 \cite{ref171} & Cross-sectional analysis, n=190 & \seqsplit{Individuals} with IPV histories assessed via SARA-V3, Canada (Indigenous vs non-Indigenous) & Structured-\seqsplit{professional}-judgment risk ratings (not only actuarial tools) showed racial \seqsplit{disparities} between Indigenous and non-Indigenous \seqsplit{participants} in risk factors and summary ratings. & Complicates\\
\\[3pt]
Lee Y 2019 \cite{srp29384426} & Randomized controlled pilot trial (online modules, baseline and 3-month follow-up) & Korean American faith leaders, \seqsplit{southeastern} USA (n=55; \seqsplit{intervention} n=27, control n=28) & The \seqsplit{intervention} group showed \seqsplit{significant} \seqsplit{improvement} in knowledge of IPV resources and attitudes against IPV compared to control, but gains in self-efficacy and prevention\slash{}\seqsplit{intervention} behaviours, while greater than control, were not \seqsplit{statistically} \seqsplit{significant}. & Complicates\\
\\[3pt]
Onukwugha FI 2025 \cite{ref108} & Mixed-methods \seqsplit{implementation} study embedded in a matched-pair cluster-controlled trial (75 providers, 20 facilities) & Antenatal care and family planning clinics, Nigeria; nurse-midwives and community health extension workers trained in WHO LIVES + ARCHES & [CORRECTION: first\_author is wrong: search lists 'Onukwugha FI' but PubMed indexes the first author as Oduenyi. Content is otherwise well supported: 'Quality assurance standards increased over time (P $\le$ 0.05)'; 'average scores of \seqsplit{approximately} 3 out of a maximum of 4'; 'providers initially perceived the \seqsplit{intervention} as additional work...lack of funds and limited staff resources were found to be key barriers'; 'pulse survey showed no \seqsplit{significant} change over time for \seqsplit{feasibility}, \seqsplit{appropriateness} and \seqsplit{acceptability}.' FIX: change first\_author to Oduenyi.] Quality-assurance readiness scores for GBV first-line response increased \seqsplit{significantly} over 9 months (P<=0.05); providers rated \seqsplit{integration} as \seqsplit{appropriate}\slash{}acceptable (\textasciitilde{}3 of 4 on average) but reported it initially felt like added work, and lack of funds\slash{}staff was a key barrier; \seqsplit{feasibility}\slash{}\seqsplit{acceptability} pulse scores showed no \seqsplit{significant} change over time despite improved facility readiness. & Complicates\\
\\[3pt]
US Preventive Services Task Force 2018 \cite{srp30357305} & guideline \slash{} \seqsplit{recommendation} statement & US women of \seqsplit{reproductive} age; older\slash{}vulnerable adults & With moderate certainty, screening plus referral to ongoing support services has moderate net benefit; explicitly states evidence does NOT support the \seqsplit{effectiveness} of brief \seqsplit{interventions} or merely providing referral \seqsplit{information} in the absence of ongoing supportive \seqsplit{intervention} components, and that benefit evidence is \seqsplit{concentrated} in pregnant\slash{}postpartum \seqsplit{populations}. & Complicates\\
\\[3pt]
Bundock Kerrie 2020 \cite{srp29683049} & Systematic review (19 studies) & Adolescent victims of dating violence, mixed \seqsplit{international} samples (English-language literature) & \seqsplit{Adolescents} \seqsplit{overwhelmingly} favored informal help sources (friends most common) over formal services; females were generally more likely than males to seek help (with \seqsplit{inconsistencies}); \seqsplit{measurement} and \seqsplit{definitions} of both violence and help-seeking varied \seqsplit{substantially} across included studies, limiting \seqsplit{comparability}. & Relevant only\\
\\[3pt]
Gutowski 2023 \cite{srp35611345} & instrument-\seqsplit{development} study, n=222 survivor-mothers, \seqsplit{exploratory} factor analysis + Rasch analysis & US survivor-mothers involved in family law \seqsplit{proceedings} & Developed a 14-item, two-subscale Legal Abuse Scale (LAS) to measure how family-court and legal processes can be misused by partners to continue coercive control post-separation. & Relevant only\\
\\[3pt]
Hameed M 2020 \cite{srp32608505} & Cochrane systematic review and meta-analysis of 33 RCTs (n=5517 women) & Women aged 16+ with recent or lifetime IPV experience, mostly high-income countries, recruited from healthcare, community, shelter\slash{}refuge settings & \seqsplit{Psychological} therapies probably reduce depression at medium-term follow-up (SMD -0.24, 95\% CI -0.47 to -0.01; 4 trials, 600 women, moderate certainty) and may reduce anxiety, but evidence is uncertain or absent for self-efficacy, PTSD, re-exposure to IPV, and safety planning; harm data are limited (only 1 of 33 trials used a validated harm-outcome scale). & Relevant only\\
\\[3pt]
Michaels-Igbokwe C 2016 \cite{srp26924488} & Economic evaluation alongside a cluster randomized controlled trial & SASA! \seqsplit{intervention} \seqsplit{communities}, Kampala, Uganda, provider \seqsplit{perspective}, 4 years of \seqsplit{programming} & Total \seqsplit{intervention} cost over 4 years was estimated at US\$553,252, averaging US\$5 per person reached annually; an estimated 1,201 cases of past-year physical IPV were averted (90\% CI 97-2,307) at a cost of US\$460 per case averted. & Relevant only\\
\\[3pt]
Schenck D 2018 \cite{ref160} & Ethical analysis \slash{} case-based policy proposal & Domino liver donor transplant candidates (pediatric case example) & [CORRECTION: Abstract quote: 'We propose a Domino Donor Advocate-based on the concept of the \seqsplit{independent} living donor advocate... We discuss the need for 2 distinct consent processes separated in time to ensure that potential domino donors... give a truly voluntary consent.' Key finding accurate. Transplant-medicine analogy again phrased as 'supports R4 no-stake assessor and P5 referral timing' --- FIX to analogy language per \seqsplit{instructions}.] Proposes a 'Domino Donor Advocate', modeled explicitly on the \seqsplit{independent} living donor advocate, plus two consent processes separated in time, to protect \seqsplit{voluntariness} when the same party stands to gain from securing consent. & Relevant only\\
\\[3pt]
Stern E 2018 \cite{ref114} & \seqsplit{Implementation}\slash{}adaptation review (\seqsplit{qualitative} evaluation and programme \seqsplit{documentation}) & \seqsplit{Indashyikirwa} programme, \seqsplit{implemented} by CARE Rwanda, RWAMREC and RWN, rural Rwanda & \seqsplit{Indashyikirwa} combined a 5-month couples curriculum, sustained community activism by trained couples, women's safe spaces for dedicated support and referral of IPV survivors, and training of opinion leaders to support an enabling \seqsplit{environment}; adaptation was guided by the SASA! fidelity brief. & Relevant only\\
\\[3pt]
World Health \seqsplit{Organization} 2013 (no PMID) & Global clinical and policy guideline & Global health-system guidance for responding to women subjected to IPV or sexual violence & \seqsplit{Establishes} WHO's first-line support principles for clinicians (the basis of the later-named LIVES model: Listen, Inquire about needs, Validate, Enhance safety, Support\slash{}refer), recommends against mandatory universal screening in the absence of adequate response systems, and sets minimum \seqsplit{requirements} (privacy, \seqsplit{confidentiality}, trained staff, referral pathways) before any clinical inquiry about violence should occur. & Relevant only\\
\\[3pt]
World Health \seqsplit{Organization} 2014 (no PMID) & Global clinical practice handbook \seqsplit{operationalizing} the 2013 WHO guidelines & Frontline health workers globally, first-line clinical response to women disclosing IPV or sexual violence & \seqsplit{Operationalizes} the LIVES first-line support protocol (Listen, Inquire, Validate, Enhance safety, Support) into a step-by-step clinical tool for frontline providers, with explicit guidance on \seqsplit{confidentiality}, \seqsplit{documentation}, safety planning, and structured referral rather than passive \seqsplit{information} provision. & Relevant only\\
\\[3pt]
\end{longtable}
